% Gauss, Werke X/1 (1917), native98, printed91.
\SourceLeaf{098}\NoSourceHead
\SourceGroup{IV}
\HistoricTitle{[ZUR BESTIMMUNG DER KLASSENANZAHL.]}
\HistoricCitation{[Aus Scheda Ae, Varia, Julius 1800, S. 36.]}
\noindent\hspace*{\parindent}Es sei $2E$ die Zahl nach der die Form der Th[eiler] recurrirt,
\[
\cos\frac{P}{4E}+i\sin\frac{P}{4E}=\cos A+i\sin A=\rho,
\]
\noindent so ist
\[
\begin{aligned}
&\frac{\surd D}{E}\{\operatorname{cotg}A\pm\operatorname{cotg}3A\pm.\,.\,.\pm\operatorname{cotg}nA\}\quad n<E\\[4pt]
={}&\frac{\surd D}{2E}\{\operatorname{cotg}A\pm\operatorname{cotg}3A\pm.\,.\,.\pm\operatorname{cotg}nA\}\quad n<2E.
\end{aligned}
\]
\ClosingRule
\HistoricCitation{[Aus Scheda Af, Mémoires de Mathematiques, Bronsuic 1801, S. 78.]}
\noindent\hspace*{\parindent}Es sei $2E$ die Zahl, nach der die Formen der Theiler recurriren
\[
\rho^{2E}=1,
\]
\noindent so ist die Zahl der quadratischen Formen
\[
\begin{aligned}
={}&\frac{\surd D}{E}\left\{\tfrac{1+\rho}{1-\rho}\pm\tfrac{1+\rho^3}{1-\rho^3}\pm.\,.\,.\pm\tfrac{1+\rho^n}{1-\rho^n}\right\}&\quad n<E\\
&&\text{(proxime)}\\[4pt]
={}&\frac{\surd D}{2E}\left\{\tfrac{1+\rho}{1-\rho}\pm\tfrac{1+\rho^3}{1-\rho^3}\pm.\,.\,.\pm\tfrac{1+\rho^n}{1-\rho^n}\right\}&\quad n<2E.
\end{aligned}
\]
\begin{center}\small BEMERKUNG.\end{center}
\begingroup\small
\noindent\hspace*{\parindent}Vergleiche hierzu die Tagebuchnotizen Nr. 114 vom 30. Oktober und Nr. 115 vom 3. Dezember 1800.
Ferner Werke II, S. 285, 286 und die zugehörigen Anmerkungen von R. \textsc{Dedekind} ebenda, S. 297,
sowie den Artikel \textbf{27} des \textsc{Bachmann}schen Aufsatzes »Über \textsc{Gauss}’ zahlentheoretische Arbeiten«.
\par\hfill\textsc{Schlesinger.}\par
\endgroup
\ClosingRule
\par\smallskip\noindent\hfill{\small 12\textsuperscript{*}}\par
\clearpage
